\documentclass[10pt]{article}
\usepackage{graphicx}
\usepackage{amssymb}

\textwidth = 6.5 in
\textheight = 9 in
\oddsidemargin = 0.0 in
\evensidemargin = 0.0 in
\topmargin = 0.0 in
\headheight = 0.0 in
\headsep = 0.0 in
\parskip = 0.1in
\parindent = 0.0in

\newenvironment{packed_enum}{
\begin{enumerate}
  \setlength{\itemsep}{1pt}
  \setlength{\parskip}{0pt}
  \setlength{\parsep}{0pt}
}{\end{enumerate}}


\newtheorem{theorem}{Theorem}
\newtheorem{corollary}[theorem]{Corollary}
\newtheorem{definition}{Definition}

\begin{document}
\pagestyle{empty}

\begin{center}
{\bf Homework Expectations and  Grading System} - Math 300 - Discrete Mathematics %- Fall 2010
\end{center}

{\bf Expectations}

This course may have a higher level of expectations for written homework than other mathematics courses.  In this course you will need to spend time writing up your solutions.  Homework assignments must be clearly written up.

Each Exercise or Proof submitted must include:
\begin{packed_enum}
\item Problem description or claim.
\item Complete solution including all important calculations.
\item Clarifying comment to help the read understand challenging parts of the solution.
\end{packed_enum}

Each submitted problem and solutions should be self contained, meaning that it can be clearly read by anyone in the class without referring to outside material. If cannot understand you work, I will ask for a rewrite and a resubmission.

{\bf Grading}

Each homework assignment is worth up to 24 points, 9 for the three exercises and 15 for the three proofs.

\underline{Exercises} are graded on the following scale:\\
3 points - problem and correct solution clearly presented\\
2 points - problem completed with near correct solution, problem and solution not clearly presented\\
1 point - problem description missing, solution incorrect, or problem partially completed\\
0 points - little to no work submitted\\

\underline{Proofs} are graded the same way as the proof portfolio according the the following system.  Points are summed and divided by 2 so that each proof is worth 5 points.

Correctness: 0-4 points\\
4 points - Claim fully demonstrated with correct logic\\
3 points - Claim demonstrated with minor logical gap and/or incorrect statements\\
2 points - Claim partially supported with some correct logic \\
0 points - Claim is not supported\\

Structure: 0-2 points\\
2 points - proof technique clearly indicated and components of proof structure easily identifiable\\
1 point - proof technique and structure clear only after multiple readings\\
0 points - proof technique is unclear and/or proof structure is incorrect\\

Completeness: 0-2 points\\
2 points - Proof features complete sentences and clear descriptions, including correct spelling, grammar, and notation\\
1 point - Proof uses sentence fragments with partial descriptions and/or proof is overly wordy and/or there are minor problems with spelling, grammar, and notation\\
0 points - Proof is difficult to understand because of terseness and/or verboseness and/or severe problems with spelling, grammar, and notation\\

Clairity: 0-2 points\\
2 points - most Math 300 student should be able to follow this proof\\
1 point - some  Math 300 students could understand this proof with careful readings\\
0 points - proof is very difficult to follow and requires multiple readings  to understand\\

For example if you grade on a proof is

$$\frac{3+1+2+1}{2}=3.5$$

This means you received 3 points for accuracy, 1 point for presentation, 2 points for Efficiency, and 1 point for Audience Appropriateness and your total score for this proof was 3.5/5.
\end{document}
